\section{Experimental Setup}
\label{sec:setup}
We evaluate Afar-prgenv on Frontier using the current Cray PE and the latest
AFAR drop available in the system module tree. The goal is to measure
compatibility with the PE toolchain and key libraries, not peak performance.

\subsection{Systems and Software}
\label{sec:systems}
Our baseline environment loads `cpe/\peversion` and `PrgEnv-amd/\prgenvversion`
alongside `rocm/\afarrocm` and `amd/\afarrocm`. Afar-prgenv modules are provided
from `/sw/crusher/ums/compilers/modulefiles/`, and the evaluation targets
`afar-prgenv/\afarversion`. This setup mirrors the site-recommended module
sequence used by the support team.

Although the AFAR drop ships a ROCm runtime, the additional `rocm/\afarrocm`
module is retained for compatibility with PE libraries that still depend on
system ROCm paths. When evaluating older dependencies, we also verify
compatibility against ROCm 6.x where needed.
The repro harness focuses on PE and AFAR modules; ROCm compatibility for
legacy dependencies is tracked as a separate gate (e.g., `libsci-acc`).

\begin{verbatim}
module purge
module reset
module load cpe/\peversion
module load rocm/\afarrocm
module load amd/\afarrocm
module load PrgEnv-amd
module load cray-mpich/9.0.1
module load craype-x86-trento
module load craype-accel-amd-gfx90a
module use /sw/crusher/ums/compilers/modulefiles/
module load afar-prgenv/\afarversion
\end{verbatim}

\subsection{Workloads and Repro Cases}
\label{sec:workloads}
We use the repository repro suite to validate compile and link correctness:
`compile-only`, `cmake-mpi-omp`, `wrapper-compile-link`, `mpi-hello`, `omp-c`,
`omp-f90`, `libsci`, `libsci-acc`, `hdf5`, `fftw`, `dsmml`, and `netcdf`. The
tests exercise the PE wrappers, MPI module files, and library metadata via
\texttt{pkg-config}, and \texttt{scripts/\allowbreak check\_ldd.sh} validates the
resulting binaries.
The CMake repro requires \texttt{cmake} in \texttt{PATH}.

The repro harness (\texttt{scripts/run\_afar\_tests.sh}) runs profiles that load each
library module both before and after `afar-prgenv`, ensuring the wrapper
re-synchronization logic works under module-order permutations.

In this study we ran:
\begin{verbatim}
scripts/run_afar_tests.sh --keep-going \
  --afar-version \afarversion \
  --afar-module afar-prgenv/\afarversion
\end{verbatim}

The default profiles exercised by the harness are summarized in
Table~\ref{tab:profiles}. Each profile extends the baseline module stack with a
library module and its associated `pkg-config` metadata, enabling focused
validation of PE libraries under AFAR.

\begin{table}[t]
  \centering
  \small
  \setlength{\tabcolsep}{3pt}
  \caption{Default repro profiles and module extensions.}
  \label{tab:profiles}
  \begin{tabular}{l l l}
    \hline
    Profile & Module(s) & Purpose \\
    \hline
    base & --- & compiler/MPI/offload sanity \\
    libsci & cray-libsci & BLAS/LAPACK links \\
    libsci-acc & cray-libsci\_acc & accelerator libs \\
    hdf5 & cray-hdf5-parallel & HDF5 parallel stack \\
    hdf5-serial & cray-hdf5 & HDF5 serial stack \\
    fftw & cray-fftw & FFTW bindings \\
    dsmml & cray-dsmml & DSMML links \\
    netcdf & cray-parallel-netcdf & PnetCDF stack \\
    \hline
  \end{tabular}
\end{table}

The harness can restrict module ordering with \texttt{--lib-order} and can
limit profiles with \texttt{--profiles}. Environment overrides such as
\texttt{AFAR\_TEST\_CPE} and \texttt{AFAR\_TEST\_MPICH} enable alternate PE
combinations, while \texttt{AFAR\_TEST\_RUN} or \texttt{AFAR\_TEST\_RUN\_MPI}
trigger runtime execution when nodes are available. These controls let support
teams focus on regression hotspots without reworking individual repro scripts.

\subsection{Metrics}
\label{sec:metrics}
The primary metric is functional readiness: all repro cases must compile and
link with the PE wrappers and AFAR drop. Secondary metrics include the
diagnostic completeness of wrapper logs, P/S/F rates per test, and the time to
validate a new drop. We also track module-order robustness by running the repro
suite under multiple module load permutations to ensure metadata
re-synchronization remains correct.
